% Input SHA256 c34ec98f0093c93effb5a37fbaeb8c8739b1d7f5ad263825e3272317a047eaec benchmark/analysis/outputs/windows-round2-intake-v1/summary.json
第二轮接收把无约束路径、参数变换与下游诊断分别核对。
108份MH保存路径均在原容差内，接受事件失配为0；但从同一无约束状态在Mac重算原尺度输出时，
60份不满足原逐位相等要求，最大绝对差为$8.88\times10^{-15}$。
这项严格跨系统检查仍记为失败，不能通过事后添加容差改称全部通过。
Windows同机核验与Mac跨系统重算回答不同问题。

CPU Pyro NUTS在9个目标上的串行/四进程技术对照中，54组保存数组逐字节相同。
对收到的18份原二进制重新计算R诊断，有限$\hat R$与tail ESS一致，
bulk ESS最大差约$1.02\times10^{-12}$，不可判定位置一致。
这不意味着重复计算后验函数也必然逐位相同，更不增加独立统计重复数。
原短链的探索不足仍然成立。

实际数组传递同样是复现契约的一部分。Windows重新生成的六份候选输入中，
噪声与导数探测数组相同，但对数均匀数组有82/49152个元素不同，最大差约$4.44\times10^{-16}$。
因此原协议正确地拒绝启动192项后续机制工作流。
已补齐原始输入文件，尚无该批执行结果；相同种子和依赖版本不能替代实际数组校验。
上述核验定位了差异层次，不证明特定编译器或数学库的因果责任。
